\documentclass[10pt,UTF8]{article}

\usepackage{geometry}

\usepackage{ctex}

\geometry{a4paper,scale=0.9}
\ctexset{today=old}

\usepackage[bookmarksnumbered=true]{hyperref}  % 生成大纲


\usepackage{times}  % 设置正文字体为 Adobe Times Roman

% \usepackage{makecell}
% \usepackage{multirow} % 表格多行合并

% \usepackage{graphicx} %插入图片的宏包
% \usepackage{subfigure}
% \usepackage{float} %设置图片浮动位置的宏包

\usepackage{amsmath}
\usepackage{amssymb}
\usepackage{mathtools}

\usepackage{physics}

% \usepackage{siunitx}
% % 关闭警告
% \ExplSyntaxOn
% \msg_redirect_name:nnn { siunitx } { physics-pkg } { none }
% \ExplSyntaxOff

%%%%%%%%%%%%%%%%%%%%%%%% tcolorbox %%%%%%%%%%%%%%%%%%%%%%%%%%
\usepackage[most]{tcolorbox}

\newenvironment{definition box}
{\begin{tcolorbox}[colback=blue!5!white,colframe=blue!75!black,parbox=false,before upper=\par,before lower=\par]}
{\end{tcolorbox}}

\newenvironment{theorem box}
{\begin{tcolorbox}[colback=yellow!5!white,colframe=yellow!75!black,parbox=false,before upper=\par,before lower=\par]}
{\end{tcolorbox}}
%%%%%%%%%%%%%%%%%%%%%%%%%%%%%%%%%%%%%%%%%%%%%%%%%%%%%%%%%%%%
            
%%%%%%%%%%%%%%%%%%%%% amsthm %%%%%%%%%%%%%%%%
\usepackage{amsthm}

\theoremstyle{definition}
\newtheorem{definition inner}{Definition}[section]
\newenvironment{definition}[1][]
{\begin{definition box}\begin{definition inner}[#1]\hfill\par}
{\end{definition inner}\end{definition box}}

\theoremstyle{plain}
\newtheorem{theorem inner}{Theorem}[section]
\newenvironment{theorem}[1][]
{\begin{theorem box}\begin{theorem inner}[#1]\hfill\par}
{\end{theorem inner}\end{theorem box}}

\theoremstyle{definition}
\newtheorem*{proof inner}{\textit{Proof}}
\renewenvironment{proof}[1][]
{\begin{proof inner}[#1]\hfill\par}
{\qed\end{proof inner}}

\theoremstyle{remark}
\newtheorem*{remark}{\textit{Remark}}
%%%%%%%%%%%%%%%%%%%%%%%%%%%%%%%%%%%%%%%%%%%%%%


\newcommand{\mycases}[2][0.8]{
    \left\{\vphantom{\scalebox{#1}{
        $\begin{matrix}#2\end{matrix}$
    }}\right.\!\!\!\!
    \begin{array}{l@{\quad}l}#2\end{array}
}

\newcommand{\ju}{\mathrm{j}}  % 虚数单位
\newcommand{\e}{\mathrm{e}{\mkern 1 mu}}


\begin{document}


\part{Linear Constant-Coefficient Differential Equations}

An $N$th-order LCCDE is given by
\begin{equation*}
    \sum_{k=0}^N a_k \dv[k]{y(t)}{t}
    =
    \sum_{k=0}^M b_k \dv[k]{x(t)}{t}
\end{equation*}

\section{Solution}

The solution to an LCCDE consists of two parts:
a \emph{homogeneous solution} plus a \emph{particular solution}.
\begin{equation*}
    y(t) = y_h(t) + y_p(t)
\end{equation*}

\subsection{Homogeneous Solution \& Natural Response}

The homogeneous solution is the solution to the homogeneous equation
\begin{equation*}
    \sum_{k=0}^N a_k \dv[k]{y(t)}{t} = 0
\end{equation*}
It has $N$ constants to be determined by the $N$ auxiliary conditions.

\paragraph{Characteristic Polynomial}
\begin{equation*}
    P(\lambda) := \sum_{k=0}^N a_k \lambda^k
\end{equation*}

\paragraph{Characteristic Equation}
\begin{equation*}
    P(\lambda) = \sum_{k=0}^N a_k \lambda^k = 0
\end{equation*}

\vspace{20pt}

\paragraph{}
If $P(\lambda)=0$ has $N$ distinct complex roots,
$\lambda_1, \lambda_2, \cdots, \lambda_N$,
then the solution is
\begin{equation*}
    y_h(t) = \sum_{k=1}^N C_k \e^{\lambda_k t}
\end{equation*}
where $C_1, C_2, \cdots, C_N$ are constants to be determined.

\paragraph{}
If $\displaystyle P(\lambda)=\prod_{k=1}^{n}(\lambda-\lambda_k)^{r_k}$,
where $\lambda_1,\lambda_2,\cdots,\lambda_n$ are distinct roots of $P(\lambda)=0$,
then the solution is
\begin{equation*}
    y_h(t) = \sum_{k=1}^n \left(
        C_{k,r_k-1}t^{r_k-1} + \cdots + C_{k,1}t + C_{k,0}
    \right) \e^{\lambda_k t}
\end{equation*}
where $C_{1,0},\cdots,C_{n,r_n-1}$ are $N$ constants to be determined.

\subsection{Particular Solution \& Forced Response}




\part{Linear Constant-Coefficient Difference Equations}


\end{document}
